\chapter{Data and Preprocessing}
\label{chap:data-preprocessing}

\section{PhysioNet 2012 Pipeline}
\label{sec:data-preprocessing:physionet-2012-pipeline}
% Section ID: C4.1.
% Evidence map: see thesis-writing/planning/chapter_evidence_map.md.

The PhysioNet/Computing in Cardiology Challenge 2012 dataset provides ICU time-series records and associated in-hospital mortality outcomes for a mortality-prediction challenge setting \cite{silva2012physionet}.  In this thesis it functions as one of the two ICU time-series sources used to exercise the proxy-state prediction and downstream causal-analysis pipeline.  The local repository does not track the raw challenge files or a verified processed-data pickle, so this section describes the implemented preprocessing contract rather than asserting a raw preprocessing cohort total.  Chapter~\ref{chap:results} separately reports the validated causal-analysis model-row count.

The causal-repository PhysioNet preprocessing implementation traverses the raw \texttt{set-a}, \texttt{set-b}, and \texttt{set-c} folders and reads one patient-record file at a time.  It removes rows with missing parameter names, skips records with at most five retained measurement rows, and removes observations with negative values, which the source code treats as missingness indicators.  It converts source times from \texttt{HH:MM} strings into elapsed minutes, renames \texttt{RecordID} to \texttt{ts\_id}, \texttt{Parameter} to \texttt{variable}, and \texttt{Value} to \texttt{value}, and reads outcome files by mapping \texttt{In-hospital\_death} to \texttt{in\_hospital\_mortality}.  After concatenating the source sets, it drops duplicate events and converts the categorical \texttt{ICUType} measurement into one-hot-style variables \texttt{ICUType\_1} through \texttt{ICUType\_4}.

The resulting causal processed artifact is serialized as three objects: \texttt{ts}, \texttt{oc}, and \texttt{ts\_ids}.  The \texttt{ts} object is a long event table with \texttt{ts\_id}, \texttt{minute}, \texttt{variable}, and \texttt{value}.  The \texttt{oc} object contains \texttt{ts\_id}, \texttt{length\_of\_stay}, \texttt{in\_hospital\_mortality}, and a source subset indicator.  The \texttt{ts\_ids} object is derived from identifiers observed in \texttt{ts}.  This artifact is the expected input to rule-based tagging, mortality-from-proxy analysis, matching, and CATE estimation in the causal repository.

\section{MIMIC-III Pipeline}
\label{sec:data-preprocessing:mimic-iii-pipeline}
% Section ID: C4.2.
% Evidence map: see thesis-writing/planning/chapter_evidence_map.md.

MIMIC-III is a critical-care electronic health-record database used here as the second ICU data source \cite{johnson2016mimiciii}.  The implemented workflow follows a clinical time-series benchmark style in that it transforms raw ICU events into stay-level time-series examples suitable for prediction, while adapting the output contract for this thesis pipeline \cite{harutyunyan_2019_mimiciii_benchmark}.  The repository does not establish a tracked raw-data snapshot, extraction date, or final processed artifact hash.

The active causal-repository MIMIC preprocessing implementation reads broad categories of MIMIC-III source tables: ICU-stay timing from \texttt{ICUSTAYS}, demographics from \texttt{PATIENTS}, bedside measurements from \texttt{CHARTEVENTS}, laboratory measurements from \texttt{LABEVENTS}, fluid outputs from \texttt{OUTPUTEVENTS}, administered inputs from \texttt{INPUTEVENTS\_CV} and \texttt{INPUTEVENTS\_MV}, and mortality/timing information from \texttt{ADMISSIONS}.  Large event tables are processed in chunks and written to temporary shards before feature rows are collected.  The implementation filters to adult ICU stays, removes chart events marked with an error flag, requires valid chart times and numeric or textual values, applies item-ID mappings and range filters for selected variables, converts selected units, and assigns events without direct ICU-stay identifiers to an ICU interval when possible.

For the causal contract, event times are converted to elapsed minutes relative to ICU admission, events are restricted to the early ICU window used by the implementation, age and gender are added as static rows at minute zero, and duplicate \((\texttt{ts\_id}, \texttt{minute}, \texttt{variable})\) events are collapsed by averaging.  The helper contract canonicalizes stay identifiers, normalizing numeric-looking identifiers such as \texttt{12345.0} to \texttt{12345}.  It defines the canonical \texttt{ts} columns as \texttt{ts\_id}, \texttt{minute}, \texttt{variable}, and \texttt{value}, defines the canonical \texttt{oc} columns as \texttt{ts\_id}, \texttt{length\_of\_stay}, \texttt{in\_hospital\_mortality}, and \texttt{subset}, and intentionally excludes internal MIMIC identifiers such as \texttt{HADM\_ID} and \texttt{SUBJECT\_ID} from the exported outcome table.

Inspection found a source-level issue that should be resolved before relying on local MIMIC preprocessing execution: the active script reads \texttt{ADMISSIONS.csv} with \texttt{DEATHTIME} for early-death filtering, while the canonical outcome helper requires \texttt{HOSPITAL\_EXPIRE\_FLAG} to construct \texttt{in\_hospital\_mortality}.  This does not change the documented target contract, but it remains a source-level reproducibility risk.

\section{STraTS Split-Aware Artifacts Versus Causal Artifacts}
\label{sec:data-preprocessing:strats-split-aware-artifacts-versus-causal-artifacts}
% Section ID: C4.3.
% Evidence map: see thesis-writing/planning/chapter_evidence_map.md.
% Source plan: F-DATAFLOW-01 in thesis-writing/planning/figure_plan.md.

The predictive STraTS repository uses a different processed-data contract from the causal repository.  Its loader expects a five-object pickle consisting of \texttt{events}, \texttt{oc}, \texttt{train\_ids}, \texttt{val\_ids}, and \texttt{test\_ids}, where the final three objects explicitly encode train, validation, and test identifiers.  The loader accepts stay identifiers named \texttt{ts\_id}, \texttt{icustay\_id}, or \texttt{ICUSTAY\_ID}, canonicalizes them to \texttt{ts\_id}, and validates the split arrays after the same normalization.  This split-aware contract is therefore not interchangeable with the causal repository's unsplit \texttt{[ts, oc, ts\_ids]} artifact.

In supervised STraTS runs, the latent-label CSV is joined to the selected split identifiers after ID normalization.  The loader filters labels to the supervised IDs, raises an error when labels are missing for any supervised \texttt{ts\_id}, raises an error for duplicate normalized labels, and uses all non-\texttt{ts\_id} columns in the latent CSV as proxy-label targets.  The mortality column in \texttt{oc} is not used as the supervised proxy-label target.  Prediction export writes \texttt{ts\_id}, one probability column per proxy state, and one thresholded binary column per proxy state; the inspected wrapper scripts export the \texttt{all} split for their prediction CSVs, although final archive-copy provenance remains incomplete.

The predictive loader also implements split-aware leakage controls.  It keeps only variables observed in the training split, derives normalization statistics from the training split, computes static-feature normalization from training rows, and validates label alignment before training.  For MIMIC-III supervised prediction it restricts events to the first 24 hours.  The pretraining path removes test data, uses training-derived variable statistics, uses a 48-hour maximum window for PhysioNet 2012, and permits MIMIC-III observations up to five days while sampling 24-hour input windows.  These controls reduce some sources of leakage, but they do not guarantee absence of leakage, nor do they resolve missing raw-data and split-provenance limitations.

The parent router can bridge contracts by reading a causal \texttt{[ts, oc, ts\_ids]} artifact and constructing a STraTS-style \texttt{[ts, oc, train\_ids, val\_ids, test\_ids]} payload through a seeded identifier shuffle.  That bridge is an orchestration convenience and does not prove how every archived final STraTS artifact was generated.

\begin{table}[htbp]
\centering
\caption{Processed artifact contracts used by the causal and STraTS repositories.}
\label{tab:data-contracts}
\small
\begin{tabular}{L{0.20\textwidth}L{0.31\textwidth}L{0.39\textwidth}}
\toprule
Property & Causal repository & STraTS repository \\
\midrule
Processed artifact & \texttt{[ts, oc, ts\_ids]} & \texttt{[events, oc, train\_ids, val\_ids, test\_ids]} \\
Split representation & Unsplit ID collection & Explicit train/validation/test IDs \\
Main use & Tagging and causal analysis & Predictive training and export \\
Identifier convention & Canonical \texttt{ts\_id} & Loader-normalized \texttt{ts\_id} \\
Proxy-label source & Separate latent CSV downstream & Separate latent CSV joined to split IDs \\
\bottomrule
\end{tabular}
\end{table}

Several provenance limitations remain.  Raw PhysioNet and MIMIC-III data are not tracked in Git.  The processed pickles used by final runs are external or absent from the audited repository.  Some run records contain machine-specific absolute paths, and final cohort counts require verified processed artifacts or manifests.  Clean-checkout reproduction therefore requires authorized data access, artifact hashes, and a documented split-generation record.
